% ACL-style paper skeleton.
%
% 1. Put acl.sty and acl_natbib.bst in this directory:
%      https://github.com/acl-org/acl-style-files
%    (or use the official Overleaf template and copy the .tex back here)
% 2. `make paper`
%
% WRITE ORDER (do not write top to bottom):
%   3 Methods -> 4 Data -> 5 Results -> 6 Analysis -> 2 Background ->
%   7 Related -> 1 Intro -> 8 Discussion -> Abstract LAST.
% Keep numbers in results/tables/*.json and \input them; never retype a number.

\documentclass[11pt]{article}
\usepackage[review]{acl}
\usepackage{times}
\usepackage{latexsym}
\usepackage[T1]{fontenc}
\usepackage[utf8]{inputenc}
\usepackage{microtype}
\usepackage{graphicx}
\usepackage{booktabs}
\usepackage{amsmath}
\usepackage{amssymb}

\title{Polarity Circuits: Is Natural-Logic Monotonicity\\Implemented in Small Language Models?}

\author{Anonymous}

\date{}

\begin{document}
\maketitle

\begin{abstract}
% Three sentences: the question, the instrument, the finding.
% If you cannot state the finding in one sentence, it is not a paper yet.
\end{abstract}

\section{Introduction}
% 1. Models pass monotonicity benchmarks. 2. Passing is not implementing.
% 3. We build a controlled instrument and test implementation causally.
% 4. Contribution bullets (3, no more).

\section{Background: The Monotonicity Calculus}
% Polarity assignment; upward/downward entailment; what a model must do to
% count as implementing it. Keep this to ~0.75 page.

\section{A Controlled Benchmark with Known Derivations}
% Generator, fragment, gold labels, human validation of a sample, difficulty knobs.

\section{Do Models Track Polarity? Behavioural Evidence}
% Controlled set + FLog + MED. Report the negative controls: if a bag-of-words
% baseline does well, say so -- it bounds the interpretation.

\section{Causal Abstraction of the Polarity Variable}
% Interchange interventions, DAS, layerwise results, null results.
% State the falsification conditions from docs/research_notes.md up front.

\section{Are KAN Probes More Faithful than Linear Probes?}
% Matched-parameter comparison on the same variable. One table, one paragraph
% of interpretation. This is a subsection, not a pillar.

\section{Related Work}

\section{Discussion and Limitations}
% Ecological validity, fragment size, single language, linear-abstraction
% boundary conditions.

\section{Conclusion}

\bibliography{refs}

\end{document}
